\scanpage{215}
\mspage{110}

Dans ce contexte général, si on part d'une ``vraie'' structure de graphe (sans boucles, i.e. $A_0\subset\mathscr P_2(S)$) séparé, il n'est pas clair qu'en \unsure{le|la} transformant par un élément de $Z$, faisant

\[
\text{donc}\quad A_\lambda=\{\{s,\lambda t\}\in\mathscr P(S)\mid s,t\in S\ \text{t.q.}\ \{s,\unsure{1|t}\}\in A_0\},
\]

on trouve encore des vrais graphes -- car il pourrait arriver qu'on ait $t=\lambda s$ pour un \deleted{$\lambda$} triple convenable $(s,t,\lambda)\in S\times S\times Z$, avec $\{s,t\}\in A_0$.

Prenons p. ex. cas $S^\vee=S$, $R={}$ diagonale, donc les jauges (satisfaisant 70 (1)) sont les $\varphi\in\unsure{\mathfrak S|\mathscr S}_S$ telles que la relation \deleted{$\unsure{(s,\varphi\in)}$} $\varphi(t)=s$ en $s,t$ soit symétrique, i.e. équivalente à $\varphi(s)=t$ i.e. à $s=\varphi^{-1}t$, ce qui signifie aussi

\[
\varphi^2=\mathrm{id}\qquad\text{i.e. }\varphi\text{ est une involution} \tag*{(71)}
\]

Les autom. \deleted{de} $(\lambda,\lambda^\vee)$ de $\Gamma=(S,S^\vee,\unsure{\mathrm{Diag}|R_{\mathrm{diag}}})$ sont les couples $(\lambda,\lambda)$ avec $\lambda\in\mathfrak S_S$, et la condition 70 (2) s'écrit

\[
\lambda\varphi\lambda=\varphi\qquad\text{soit}\quad\varphi(\lambda)=\lambda^{-1} \tag*{(72)}
\]

i.e. l'involution ``renverse'' $\lambda$. Si donc on se donne un sous-groupe commutatif $Z$ de $\mathfrak S_S$, alors $\Pi_Z$ est formé des involutions \inserted{$\varphi\in\mathfrak S_S$} qui \deleted{\unsure{opè}} normalisent $Z$, en y induisant l'autom. \deleted{$\unsure{Z}$} $\lambda\longmapsto\lambda^{-1}$. Si p. ex. $S$ est \inserted{l'ens. sous-jacent à} un torseur $L$ sous $Z$, \inserted{(donc $Z$ est son propre centralisateur dans $\mathfrak S_S$),} \deleted{ce sont les $\varphi$ en question s'identifient} \deleted{avec involutions} alors choisissant une origine dans $L$ d'où $L\simeq Z$, ce sont les involutions de la forme

\[
\varphi_\lambda(z)=\lambda-z\qquad \lambda\in Z,\ z\in Z,
\]

Ici la relation \deleted{$\unsure{z\in}\ s=\varphi_a(s)\ \text{i.e. }z=a-z$}

% archive sheet number: 214
